\documentclass[10pt,a4paper]{amsart}
\usepackage[margin=19mm]{geometry}
\usepackage{amsmath,amssymb,mathtools}
\usepackage[scr=rsfs,cal=euler]{mathalfa}
\usepackage{xfrac}
\usepackage[unicode,hidelinks,bookmarks=false]{hyperref}
\newcommand{\ie}{i.e.\ }
\newcommand{\cf}{cf.\ }
\newcommand{\resp}{respectively\ }
\newcommand{\Subsection}{\subsection}
\providecommand{\nobreakdash}{\nobreak\hbox{-}}
\setlength{\emergencystretch}{2em}
\multlinegap0pt
\usepackage{bm}
\usepackage{bbm}
\usepackage{dsfont}
\usepackage{xspace}
\usepackage{cancel}
\usepackage{mathtools}
\usepackage{amsthm}
\makeatletter
\@ifundefined{theorem}{\newtheorem{theorem}{Theorem}}{}
\@ifundefined{lemma}{\newtheorem{lemma}[theorem]{Lemma}}{}
\makeatother
\DeclareMathOperator{\supp}{supp}
\title[Continuous solutions of the complex Kac--Bernstein functiona]{Continuous solutions of the complex Kac--Bernstein functional equation on the integers and the real numbers}
\author{Takashi Satomi}
\date{Source: arXiv:2607.24033v1 (CC BY 4.0)}
\begin{document}
\maketitle

\noindent\emph{Source and licence.} Continuous solutions of the complex Kac--Bernstein functional equation on the integers and the real numbers. Version arXiv:2607.24033v1 (CC BY 4.0), distributed under the Creative Commons Attribution 4.0 International licence, as recorded in the arXiv licence metadata for this exact version and shown on the abstract page https://arxiv.org/abs/2607.24033v1. Full rights notice: licenses/arxiv-2607.24033-CC-BY-4.0.txt.\par\smallskip

\noindent\textit{Editorial scope.} Counted unit: the Lemma whose source label is \texttt{lem:supp-large-Gauss} in the article (source0.tex lines 1168--1184, source label \texttt{lem:supp-large-Gauss}), delivered together with the article's own complete original proof of it (source0.tex lines 1186--1299, one \texttt{proof} environment of 114 lines). No mathematics is written, repaired or re-derived: statement and proof are the article's own text line for line. Required context retained uncounted: eq:complex-Kac-Bernstein (lines 78--81); thm:Z-characterize (lines 95--170); item:Z-characterize-six-point (lines 129--132); eq:Z-characterize-six-point-equation (lines 142--144); item:Z-characterize-Gauss (lines 142--144); lem:six-point-equation (lines 782--790); lem:supp-arithmetic-progression (lines 806--836); eq:supp-arithmetic-progression-S-define (lines 810--813); item:supp-arithmetic-progression-supp (lines 817--834); eq:supp-arithmetic-progression-supp-include (lines 821--824); item:supp-arithmetic-progression-main (lines 821--824); eq:supp-all-Gauss-represent (lines 912--917); lem:supp-all (lines 1001--1055); item:supp-all-Kac-Bernstein (lines 1006--1053); item:supp-all-Gauss (lines 1010--1017); lem:scaling (lines 1128--1144); eq:scaling-main (lines 1138--1141). Every definition, display and label of those lines is retained unchanged. Omitted from this package and not redistributed: 1--77, 82--94, 133--141, 145--781, 791--805, 814--816, 825--911, 918--1000, 1018--1127, 1142--1167, 1185--1185, 1300--1803. Nothing inside a retained excerpt is omitted or altered: every retained body is delivered exactly as the article's lines are written, so no omission receipt of any kind accompanies this package. Citations: the retained excerpts carry no citation command at all, so no bibliography entry of the article is transcribed and no reference is redistributed. Reference adaptation: none was needed and none was made; every reference inside the delivered unit resolves inside it, so no unresolved reference or citation is produced. Printed slips kept literal: no printed word, symbol, hypothesis or number of the counted statement or of its proof was altered. Layout and class: the article's own class is \texttt{article} with packages including geometry, amssymb, amsthm, amsmath, mathtools, hyperref, biblatex, enumitem, multirow; the wrapper is the lane's pinned \texttt{amsart} editorial wrapper at 10pt on A4, which restates the theorem-like environments the retained text needs and adds the article's own macro definitions that the retained text needs (\texttt{\textbackslash supp}), with extra wrapper packages (bm, bbm, dsfont, xspace, cancel, mathtools). The delivered numbering is the unit's own. Assets: no retained span contains a figure environment, a graphics-inclusion command, a PDF-page inclusion, a table or a TikZ picture; no image file is copied into this package, the package declares no asset dependency and no image file is redistributed with it. Licence: version arXiv:2607.24033v1 of this article is distributed under the Creative Commons Attribution 4.0 International licence, as recorded in the arXiv licence metadata for this exact version and shown on the abstract page https://arxiv.org/abs/2607.24033v1. A full rights notice is delivered at licenses/arxiv-2607.24033-CC-BY-4.0.txt. Characters: every retained excerpt is pure ASCII, so the delivered engine, XeLaTeX with its default Latin Modern text font, reports no missing character.
\begin{align}
  & \phantom{ {} = {} }  f_1 ( u + v ) f_2 ( u - v ) f_1 ( u' + v' ) f_2 ( u' - v' ) \\
  & = f_1 ( u + v' ) f_2 ( u - v' ) f_1 ( u' + v ) f_2 ( u' - v ) \label{eq:complex-Kac-Bernstein}
\end{align}

\begin{theorem}
    \label{thm:Z-characterize}

The functions $ f_1 , f_2 \colon \mathbb{ Z } \to \mathbb{ C } $ satisfy \eqref{eq:complex-Kac-Bernstein} for any $ u , u' , v , v' \in \mathbb{ Z } $ if and only if at least one of the following conditions holds:

\begin{enumerate}
  \item \label{item:Z-characterize-f1-zero}
  
  One has $ f_1 \equiv 0 $.

  \item \label{item:Z-characterize-f2-zero}
  
  One has $ f_2 \equiv 0 $.

\item \label{item:Z-characterize-f1-one-point}

There exists $ a \in \mathbb{ Z } $ such that
\begin{align}
    \supp f_1 = \{ a \} , & & 
    \# ( \supp f_2 \cap ( a + 2 \mathbb{ Z } ) ) = 1,
\end{align}
where $ \supp f \coloneqq \{ u \in \mathbb{ Z } \mid f ( u ) \neq 0 \} $ and $ \# S $ denotes the number of elements of $ S $.

\item \label{item:Z-characterize-f2-one-point}

There exists $ a \in \mathbb{ Z } $ such that
\begin{align}
    \supp f_2 = \{ a \} , & &
    \# ( \supp f_1 \cap ( a + 2 \mathbb{ Z } ) ) = 1.
\end{align}

\item \label{item:Z-characterize-independent}

There exists $ a \in \mathbb{ Z } $ such that
\begin{align}
  \supp f_1 \subset a + 2 \mathbb{ Z }, & &
  \supp f_2 \subset a + 1 + 2 \mathbb{ Z }.
\end{align}

\item \label{item:Z-characterize-six-point}

There exist $ a_1 , a_2 , k \in \mathbb{ Z } $ with $ k , a_1 + a_2 \in 1 + 2 \mathbb{ Z } $ such that
\begin{align}
    \supp f_1 \subset \{ a_1 , a_1 \pm k \}, &  &
\supp f_2 \subset \{ a_2 , a_2 \pm k \} \label{eq:Z-characterize-six-point-support}
\end{align}
and additionally
\begin{align}
    f_1 ( a_1 )^2 f_2 ( a_2 - k ) f_2 ( a_2 + k ) = f_2 ( a_2 )^2 f_1 ( a_1 - k ) f_1 ( a_1 + k ). \label{eq:Z-characterize-six-point-equation}
\end{align}

\item \label{item:Z-characterize-Gauss}

There exist $c_1 , c_2 , c_3 , c_4 , c_5 , c_6 \in \mathbb{ C } $, $ k = 1 , 3 , 5 , \cdots $, and $ a_1 , a_2 \in \mathbb{ Z } $ with $ a_1 + a_2 \in 2 \mathbb{ Z } $ such that
\begin{align}
  f_1 ( u )
  & =
  \left\{
  \begin{aligned}
    & e^{ c_1 u^2 + c_2 ( - 1 )^u + c_3 u + c_4 } & & \text{ if $ u \in a_1 + k \mathbb{ Z } $} \\
    & 0 & & \text{ if $ u \notin a_1 + k \mathbb{ Z } $ }
  \end{aligned}
  \right. , \\
  f_2 ( u )
  & =
  \left\{
  \begin{aligned}
    & e^{ c_1 u^2 - c_2 ( - 1 )^u + c_5 u + c_6 } & & \text{ if $ u \in a_2 + k \mathbb{ Z } $} \\
    & 0 & & \text{ if $ u \notin a_2 + k \mathbb{ Z } $ }
  \end{aligned}
  \right. .
\end{align}

\end{enumerate}

\end{theorem}

\begin{lemma}
    \label{lem:six-point-equation}

If functions $ f_1 , f_2 \colon \mathbb{ Z } \to \mathbb{ C } $ satisfy \eqref{eq:complex-Kac-Bernstein} for any $ u , u' , v , v' \in \mathbb{ Z } $, then \eqref{eq:Z-characterize-six-point-equation} holds for any $ a_1 , a_2 , k \in \mathbb{ Z } $ with
\begin{align}
   a_1 + a_2 + k \in 2 \mathbb{ Z }. \label{eq:six-point-equation-assume}
\end{align}

\end{lemma}

\begin{lemma}
   \label{lem:supp-arithmetic-progression}

Let $ f_1 $ and $ f_2 $ be as in Lemma \ref{lem:six-point-equation}, and
\begin{align}
    S_1 ( a , j ) \coloneqq \supp f_1 \cap ( a + j \mathbb{ Z } ), & &
    S_2 ( a , j ) \coloneqq \supp f_2 \cap ( a + j \mathbb{ Z } ) \label{eq:supp-arithmetic-progression-S-define}
\end{align}
for $ a , j \in \mathbb{ Z } $.
Suppose $ a_2 , k \in \mathbb{ Z } $ satisfy $ a_2 \pm k \in \supp f_2 $ and let $ a_1 \in S_1 ( a_2 + k , 2 ) $.

\begin{enumerate}
    \item \label{item:supp-arithmetic-progression-supp}

    One has
    \begin{align}
        a_1 \pm k \in S_1 ( a_2 , 2 ), & &
        a_2 \in S_2 ( a_1 + k , 2 ). \label{eq:supp-arithmetic-progression-supp-include}
    \end{align}

\item \label{item:supp-arithmetic-progression-main}

Furthermore, if $ k \in 2 \mathbb{ Z } $ holds, then one has
\begin{align}
    a_1 + ( k / 2 ) \mathbb{ Z } = S_1 ( a_1 , k / 2 ), & &
    a_2 + ( k / 2 ) \mathbb{ Z } = S_2 ( a_2 , k / 2 ).
\end{align}

\end{enumerate}

\end{lemma}

\begin{align}
  f_1 ( u )
  = e^{ c_1 u^2 + c_2 ( - 1 )^u + c_3 u + c_4 }, & &
  f_2 ( u )
  = e^{ c_1 u^2 - c_2 ( - 1 )^u + c_5 u + c_6 } \label{eq:supp-all-Gauss-represent}
\end{align}

\begin{lemma}
  \label{lem:supp-all}

For any $ q_1 , q_2 , q_3 , q_4 , q_5 , q_6 \in \mathbb{ C } $, the following conditions on $ f_1 , f_2 \colon \mathbb{ Z } \to \mathbb{ C } $ are equivalent.

\begin{enumerate}
    \item \label{item:supp-all-Kac-Bernstein}

    One has \eqref{eq:complex-Kac-Bernstein} for any $ u , u' , v , v' \in \mathbb{ Z } $ and
    \begin{align}
    f_1 ( - 1 ) = e^{ q_1 } , & &
    f_1 ( 0 ) = e^{ q_2 } , & &
    f_1 ( 1 ) = e^{ q_3 } , & &
    f_2 ( - 1 ) = e^{ q_4 } , & &
    f_2 ( 0 ) = e^{ q_5 } , & &
    f_2 ( 1 ) = e^{ q_6 }. \label{eq:supp-all-Kac-Bernstein-exponent}
\end{align}

\item \label{item:supp-all-Gauss}

One has \eqref{eq:supp-all-Gauss-represent} for any $ u \in \mathbb{ Z } $ when $ c_1 , c_2 , c_3 , c_4 , c_5 , c_6 \in \mathbb{ C } $ are defined as
\begin{align}
    P \coloneqq \frac{ 1 }{ 8 }
    \begin{pmatrix}
    2 & - 4 & 2 & 2 & - 4 & 2 \\
    - 1 & 2 & - 1 & 1 & - 2 & 1 \\
    - 4 & 0 & 4 & 0 & 0 & 0 \\
    1 & 6 & 1 & - 1 & 2 & - 1 \\
    0 & 0 & 0 & - 4 & 0 & 4 \\
    - 1 & 2 & - 1 & 1 & 6 & 1
  \end{pmatrix}, & &
  \begin{pmatrix}
    c_1 \\
    c_2 \\
    c_3 \\
    c_4 \\
    c_5 \\
    c_6
  \end{pmatrix}
  \coloneqq P
  \begin{pmatrix}
    q_1 \\
    q_2 \\
    q_3 \\
    q_4 \\
    q_5 \\
    q_6
  \end{pmatrix}
  .
  \label{eq:supp-all-Gauss-const-def}
\end{align}
    
\end{enumerate}

\end{lemma}

\begin{lemma}
    \label{lem:scaling}

Suppose $ a_1 , a_2 , j \in \mathbb{ Z } $ satisfy $ a_1 + a_2 \in 2 \mathbb{ Z } $, and let $ f_1 $ and $ f_2 $ be as in Lemma \ref{lem:six-point-equation}.
Then two functions $ p_1 , p_2 \colon \mathbb{ Z } \to \mathbb{ C } $ defined as 
\begin{align}
    p_1 ( u ) \coloneqq f_1 ( j u + a_1 ), & &
    p_2 ( u ) \coloneqq f_2 ( j u + a_2 )
\end{align}
satisfy the complex Kac--Bernstein functional equation 
\begin{align}
  & \phantom{ {} = {} }  p_1 ( u + v ) p_2 ( u - v ) p_1 ( u' + v' ) p_2 ( u' - v' ) \\
  & = p_1 ( u + v' ) p_2 ( u - v' ) p_1 ( u' + v ) p_2 ( u' - v ) \label{eq:scaling-main}
\end{align}
for any $ u , u' , v , v' \in \mathbb{ Z } $.
    
\end{lemma}

\begin{lemma}
    \label{lem:supp-large-Gauss}

Let $ f_1$, $ f_2 $, $ S_1$, and $ S_2 $ be as in Lemma \ref{lem:supp-arithmetic-progression}, and $ a_1 \in \supp f_1 $.

\begin{enumerate}
    \item \label{item:supp-large-Gauss-large}

If $ \# S_2 ( a_1 , 2 ) \geq 3 $, then the condition \ref{item:Z-characterize-Gauss} of Theorem \ref{thm:Z-characterize} holds.

\item \label{item:supp-large-Gauss-number-2}

If $ \# S_2 ( a_1 , 2 ) = 2 $, then the condition \ref{item:Z-characterize-six-point} of Theorem \ref{thm:Z-characterize} holds.
    
\end{enumerate}

\end{lemma}

\begin{proof}

\begin{enumerate}
    \item 

    First, we prove that there exist $ a_2 \in \supp f_2 $ and $ j_0 \in 1 + 2 \mathbb{ Z } $ such that
    \begin{align}
        a_1 + j_0 \mathbb{ Z } = \supp f_1, & &
        a_2 + j_0 \mathbb{ Z } = \supp f_2. \label{eq:supp-large-Gauss-large-proof-supp}
    \end{align}
    Since we have
    \begin{align}
        \# S_2 ( a_1 , 4 ) + \# S_2 ( a_1 + 2 , 4 ) = \# S_2 ( a_1 , 2 ) \geq 3,
    \end{align}
    at least one of $ \# S_2 ( a_1 , 4 ) \geq 2 $ and $ \# S_2 ( a_1 + 2 , 4 ) \geq 2 $ holds.
Thus, there exist $ a_2 , j \in \mathbb{ Z } $ with $ a_2 \pm 2 j \in \supp f_2 $ such that $ a_1 \in S_1 ( a_2 \pm 2 j , 2 ) $ holds.
We have
\begin{align}
    a_1 + j \mathbb{ Z } = S_1 ( a_1 , j ) , & &
    a_2 + j \mathbb{ Z } = S_2 ( a_2 , j )
\end{align}
by Lemma \ref{lem:supp-arithmetic-progression} \ref{item:supp-arithmetic-progression-main} and hence
\begin{align}
    j_0 
    \coloneqq \min \left\{ j\in\mathbb Z_{\geq1} \; \middle| \;
  \begin{aligned}
  & \text{there exist } b_1,b_2\in\mathbb Z
  \text{ with } b_1+b_2\in2\mathbb Z,\\
  & b_1+j\mathbb Z=S_1(b_1,j), \quad
  b_2+j\mathbb Z=S_2(b_2,j)
  \end{aligned}
\right\} \label{eq:supp-large-Gauss-large-proof-j0-minimum}
\end{align}
can be defined.
We have $ j_0 \in 1 + 2 \mathbb{ Z } $ by Lemma \ref{lem:supp-arithmetic-progression} \ref{item:supp-arithmetic-progression-main} and the minimality of $ j_0 $.
There exist $ b_1 , b_2 \in \mathbb{ Z } $ such that
\begin{align}
    b_1 + j_0 \mathbb{ Z } = S_1 ( b_1 , j_0 ), & &
    b_2 + j_0 \mathbb{ Z } = S_2 ( b_2 , j_0 ) \label{eq:supp-large-Gauss-large-proof-b1-b2-define}
\end{align}
by \eqref{eq:supp-large-Gauss-large-proof-j0-minimum}.

Next, we prove $ b_1 + j_0 \mathbb{ Z } = \supp f_1 $.
Since
\begin{align}
    b_1 + j_0 \mathbb{ Z } = S_1 ( b_1 , j_0 ) \subset \supp f_1
\end{align}
holds by \eqref{eq:supp-arithmetic-progression-S-define}, it suffices to show $ b_1' \in b_1 + j_0 \mathbb{ Z } $ for any $ b_1' \in \supp f_1 $.
We may assume
\begin{align}
    b_1 + b_1' \in 2 \mathbb{ Z } , & &
    b_1 \leq b_1' < b_1 + 2 j_0
\end{align}
by replacing $ b_1 $ if necessary.
By $ b_1 + b_1' \in 2 \mathbb{ Z } $ and $ j_0 \in 1 + 2 \mathbb{ Z } $, there exists $ j' \in \mathbb{ Z } $ such that either $ b_1' - b_1 = 4 j' $ or $ b_1 + 2 j_0 - b_1' = 4 j' $ holds.
Then we have $ 0 \leq j' < j_0 $.
Since $ S_2 ( b_1' , 2 ) \cap \{ b_2 , b_2 + j_0 \} \neq \emptyset $ holds by \eqref{eq:supp-large-Gauss-large-proof-b1-b2-define} and $ j_0 \in 1 + 2 \mathbb{ Z } $, we have $ j' = 0 $ by Lemma \ref{lem:supp-arithmetic-progression} \ref{item:supp-arithmetic-progression-main} and the minimality of $ j_0 $.
Thus, we have $ b_1' = b_1 \in b_1 + j_0 \mathbb{ Z } $ and hence $ b_1 + j_0 \mathbb{ Z } = \supp f_1 $.
Similarly, we obtain $ a_2 + j_0 \mathbb{ Z } = \supp f_2 $ and hence \eqref{eq:supp-large-Gauss-large-proof-supp}.

Since $ j_0 \in 1 + 2 \mathbb{ Z } $ holds, we may assume $ a_1 + a_2 \in 2 \mathbb{ Z } $ by replacing $ a_2 $ with $ a_2 + j_0 $ if necessary.
We define functions $ p_1 , p_2 \colon \mathbb{ Z } \to \mathbb{ C } $ by
\begin{align}
p_1 ( u ) \coloneqq f_1 ( a_1 + u j_0 ), & &
p_2 ( u ) \coloneqq f_2 ( a_2 + u j_0 ).
\end{align}
By Lemma \ref{lem:scaling}, the functions $ p_1 $ and $ p_2 $ satisfy \eqref{eq:scaling-main}.
Since \eqref{eq:supp-large-Gauss-large-proof-supp} implies that
\begin{align}
\supp p_1 = \supp p_2 = \mathbb{ Z },
\end{align}
Lemma \ref{lem:supp-all} $ \text{ \ref{item:supp-all-Kac-Bernstein} } \Longrightarrow \text{ \ref{item:supp-all-Gauss} } $ shows that the condition \ref{item:Z-characterize-Gauss} of Theorem \ref{thm:Z-characterize} holds.

\item

There exist $ a_2 \in \mathbb{ Z } $ and $ k \in \mathbb{ Z } \setminus \{ 0 \} $ such that $ S_2 ( a_1 , 2 ) = \{ a_2 \pm k \} $ holds by $ \# S_2 ( a_1 , 2 ) = 2 $.
Then we have $ a_1 + a_2 + k \in 2 \mathbb{ Z } $ and hence \eqref{eq:Z-characterize-six-point-equation} by Lemma \ref{lem:six-point-equation}.
In addition, we get
\begin{align}
    a_1 \in S_1 ( a_2 + k , 2 ), & &\{ a_2 \pm k \} = S_2 ( a_1 , 2 ) \subset \supp f_2
\end{align}
and hence \eqref{eq:supp-arithmetic-progression-supp-include} by Lemma \ref{lem:supp-arithmetic-progression} \ref{item:supp-arithmetic-progression-supp}.
Since
\begin{align}
    a_2 \notin \{ a_2 \pm k \} = S_2 ( a_1 , 2 )
\end{align}
holds, we have $ a_1 + a_2 , k \in 1 + 2 \mathbb{ Z } $ by \eqref{eq:supp-arithmetic-progression-supp-include}.

We now prove $ \{ a_2 , a_2 \pm k \} = \supp f_2 $.
We have $ \{ a_2 , a_2 \pm k \} \subset \supp f_2 $ by $ S_2 ( a_1 , 2 ) = \{ a_2 \pm k \} $ and \eqref{eq:supp-arithmetic-progression-supp-include}.
Thus, it suffices to show $ a_2 = a_2' $ for any $ a_2' \in \supp f_2 $ with $ a_2' \neq a_2 \pm k $.
Since $ a_2' \in S_2 ( a_1 + k , 2 ) $ holds by $ k \in 1 + 2 \mathbb{ Z } $ and 
\begin{align}
    a_2' \notin \{ a_2 \pm k \} = S_2 ( a_1 , 2 ),
\end{align}
we have
\begin{align}
    a_2' \pm k \in S_2 ( a_1 , 2 ) = \{ a_2 \pm k \}
\end{align}
by \eqref{eq:supp-arithmetic-progression-supp-include} and Lemma \ref{lem:supp-arithmetic-progression} \ref{item:supp-arithmetic-progression-supp}.
Thus, we obtain $ a_2 = a_2' $.

Finally, it suffices to show $ \# S_1 ( a_2 , 2 ) = 2 $, since $ \{ a_1 , a_1 \pm k \} = \supp f_1 $ follows by a similar argument.
The condition \ref{item:Z-characterize-Gauss} of Theorem \ref{thm:Z-characterize} does not hold by $ \# S_2 ( a_1 , 2 ) = 2 $ and hence we have $ \# S_1 ( a_2 , 2 ) \leq 2 $ by Lemma \ref{lem:supp-large-Gauss} \ref{item:supp-large-Gauss-large}.
Since
\begin{align}
   \# S_1 ( a_2 , 2 ) \geq \# \{ a_1 \pm k \} = 2 
\end{align}
holds by $ S_2 ( a_1 , 2 ) = \{ a_2 \pm k \} $, we obtain $ \# S_1 ( a_2 , 2 ) = 2 $.
\qedhere

\end{enumerate}

\end{proof}

\end{document}
